% ============================================================================
%  Préambule partagé. Compatible pdfLaTeX (pas de fontspec), donc compilable
%  tel quel sur Overleaf, MiKTeX ou TeX Live.
%  Reprend la charte du rapport Genow : géométrie, couleurs, titres compacts,
%  en-têtes fancyhdr, listings, tableaux booktabs.
% ============================================================================
\usepackage[top=2.5cm, bottom=2.5cm, left=2.5cm, right=2.5cm, headheight=15pt]{geometry}
\usepackage[utf8]{inputenc}
\usepackage[T1]{fontenc}
\usepackage[french]{babel}
\usepackage{lmodern}
\usepackage{microtype}

\usepackage{graphicx}
\graphicspath{{figures/}}
\usepackage{setspace}
\usepackage{fancyhdr}
\usepackage{amsmath,amssymb,amsthm}
\usepackage{bm}
\usepackage{booktabs}
\usepackage{array}
\usepackage{tabularx}
\usepackage{longtable}
\usepackage{multirow}
\usepackage{makecell}
\usepackage{float}
\usepackage{caption}
\usepackage{subcaption}
\usepackage{upgreek}
\usepackage{siunitx}
\sisetup{detect-all, range-phrase=--, per-mode=symbol, output-decimal-marker={,}}
\DeclareSIPrefix{\micro}{\ensuremath{\upmu}}{-6}
\usepackage{textcomp}
\usepackage{pifont}
\usepackage[dvipsnames,table]{xcolor}
\usepackage{listings}
\usepackage{tikz}
\usetikzlibrary{arrows.meta,positioning,shapes.geometric,shapes.misc,fit,calc,backgrounds}
\usepackage[most]{tcolorbox}
\usepackage{enumitem}
\setlist{itemsep=2pt, topsep=4pt}
\usepackage{url}
\usepackage[hidelinks,bookmarksnumbered,hypertexnames=false]{hyperref}

\setstretch{1.2}
\setlength{\parskip}{0.3em}
\setlength{\emergencystretch}{3em}

% Titres de chapitre compacts : numéro et titre sur un seul bloc.
\usepackage{titlesec}
\titleformat{\chapter}[hang]{\normalfont\LARGE\bfseries\color{insatblue}}{\thechapter}{0.6em}{}[\vspace{2pt}\titlerule]
\titlespacing*{\chapter}{0pt}{-10pt}{14pt}
\titlespacing*{\section}{0pt}{10pt plus 2pt minus 2pt}{5pt plus 1pt}
\titlespacing*{\subsection}{0pt}{8pt plus 2pt minus 2pt}{3pt plus 1pt}
\captionsetup{font=small, labelfont=bf, skip=5pt}
\captionsetup[sub]{font=footnotesize}
\setlength{\textfloatsep}{12pt plus 2pt minus 3pt}
\setlength{\floatsep}{10pt plus 2pt minus 2pt}
\setlength{\intextsep}{10pt plus 2pt minus 2pt}

% Permet la coupure des identifiants monospace longs (power_normalize, ...).
\makeatletter
\let\tf@origtexttt\texttt
\renewcommand{\texttt}[1]{\tf@origtexttt{\tf@breakable{#1}}}
\newcommand{\tf@breakable}[1]{\begingroup\def\_{\textunderscore\allowbreak}#1\endgroup}
\makeatother

% ------------------------------------------------------------ éléments de base
\newcommand{\cmark}{\ding{51}}
\newcommand{\xmark}{\ding{55}}
\theoremstyle{definition}
\newtheorem{definition}{Définition}[chapter]

\pagestyle{fancy}
\fancyhf{}
\fancyhead[R]{\small\nouppercase\rightmark}
\fancyfoot[R]{\textbf{\thepage}}
\renewcommand{\headrulewidth}{0.4pt}
\fancypagestyle{plain}{\fancyhf{}\fancyfoot[R]{\textbf{\thepage}}\renewcommand{\headrulewidth}{0pt}}

% -------------------------------------------------------------------- couleurs
\definecolor{insatblue}{HTML}{1F4E9A}
\definecolor{codebg}{HTML}{F7F7F5}
\definecolor{codecomment}{HTML}{6A737D}
\definecolor{codekw}{HTML}{1F4E9A}
\definecolor{codestr}{HTML}{A33B00}
\definecolor{keyfill}{HTML}{EAF1FB}
\definecolor{keyframe}{HTML}{1F4E9A}
\definecolor{blk1}{HTML}{DCEAF9}
\definecolor{blk2}{HTML}{FDE5D8}
\definecolor{blk3}{HTML}{D7F2E7}
\definecolor{blk4}{HTML}{ECEAE4}

% -------------------------------------------------------------------- listings
\lstset{
  basicstyle=\ttfamily\footnotesize,
  backgroundcolor=\color{codebg},
  frame=single, rulecolor=\color{black!20},
  breaklines=true, breakatwhitespace=false,
  numbers=left, numberstyle=\tiny\color{codecomment}, numbersep=6pt,
  keywordstyle=\color{codekw}\bfseries,
  commentstyle=\color{codecomment}\itshape,
  stringstyle=\color{codestr},
  showstringspaces=false, tabsize=2, columns=fullflexible, keepspaces=true,
  captionpos=b, xleftmargin=1.5em, framexleftmargin=1.5em,
  literate={é}{{\'e}}1 {è}{{\`e}}1 {ê}{{\^e}}1 {à}{{\`a}}1 {ç}{{\c c}}1 {ù}{{\`u}}1
           {µ}{{$\mu$}}1 {°}{{\textdegree}}1 {→}{{$\rightarrow$}}1 {—}{{--}}1 {–}{{-}}1 {×}{{$\times$}}1
}
\renewcommand{\lstlistlistingname}{Liste des extraits de code}

% Encadré pour un résultat ou une idée clé.
\newtcolorbox{keybox}[1][]{colback=keyfill, colframe=keyframe, fonttitle=\bfseries\small,
  title={#1}, boxrule=0.6pt, arc=2pt, left=6pt, right=6pt, top=4pt, bottom=4pt, fontupper=\small}

% Référence à un fichier source du projet.
\newcommand{\srcfile}[1]{\texttt{\detokenize{#1}}}
